\documentclass[11pt,a4paper]{article}
\usepackage[margin=22mm,headheight=15pt]{geometry}
\usepackage[no-math]{fontspec}
\setmainfont{NotoSerifCJKsc-Regular.otf}[Path=assets/,BoldFont=NotoSerifCJKsc-Regular.otf,BoldFeatures={FakeBold=1.7},ItalicFont=NotoSerifCJKsc-Regular.otf,ItalicFeatures={FakeSlant=.12}]
\setsansfont{NotoSerifCJKsc-Regular.otf}[Path=assets/,BoldFont=NotoSerifCJKsc-Regular.otf,BoldFeatures={FakeBold=1.7}]
\setmonofont{DejaVu Sans Mono}[Scale=.85]
\XeTeXlinebreaklocale "zh"
\XeTeXlinebreakskip=0pt plus 1pt
\usepackage{amsmath,amssymb,mathtools,mathrsfs}
\usepackage{booktabs,longtable,array,calc}
\usepackage{graphicx,xcolor,fancyhdr,enumitem,needspace}
\usepackage{float}
\floatplacement{figure}{H}
\setkeys{Gin}{keepaspectratio}
\usepackage[unicode,breaklinks]{hyperref}
\hypersetup{colorlinks=true,linkcolor=black,urlcolor=blue!55!black,pdftitle={Research MAX-STRUCT：非线性 Maxwell 逆散射的结构发现},pdfauthor={Research MAX-STRUCT}}
\urlstyle{same}
\usepackage{xurl}
\setlength{\parindent}{0pt}
\setlength{\parskip}{5.5pt}
\setlength{\emergencystretch}{3em}
\widowpenalty=10000
\clubpenalty=10000
\displaywidowpenalty=10000
\linespread{1.18}
\setcounter{tocdepth}{2}
\setcounter{secnumdepth}{0}
\renewcommand{\contentsname}{目录}
\renewcommand{\figurename}{图}
\setlist{nosep,leftmargin=1.7em,topsep=4pt}
\newcommand{\tightlist}{\setlength{\itemsep}{0pt}\setlength{\parskip}{0pt}}
\providecommand{\pandocbounded}[1]{#1}
\providecommand{\real}{\operatorname{Re}}
\providecommand{\imag}{\operatorname{Im}}
\pagestyle{fancy}
\fancyhf{}
\fancyhead[L]{\small Research MAX-STRUCT}
\fancyhead[R]{\small 理论研究 · 2026-11-02}
\fancyfoot[C]{\small\thepage}
\renewcommand{\headrulewidth}{0.3pt}
\renewcommand{\arraystretch}{1.18}
\setlength{\LTpre}{7pt}
\setlength{\LTpost}{7pt}
\begin{document}
\begin{titlepage}
\vspace*{14mm}
{\Large 非线性 Maxwell 逆散射}\par
\vspace{8mm}
{\Huge\bfseries Research MAX-STRUCT}\par
\vspace{8mm}
{\LARGE 从物理约束、可观测商与全波谱序\par 发现新的数学问题}\par
\vspace{12mm}
{\large 超越局部反馈的结构研究与实验交接}\par
\vspace{14mm}
\begin{minipage}{.92\textwidth}
本报告的主张不是已经得到一个超过成熟 CIE-I 的反演器，而是确定三条具有独立物理内容、可以被严格证明或小实验推翻的研究主线：固定磁场约束空间；有限通道全波单调性；有限照明下的同场材料族。\par
报告明确区分经典结果、仓库既有成果、本轮条件推导、待证假说及实际运行的微型验证。
\end{minipage}
\vfill
{\large 研究日期：2026 年 11 月 2 日}\par
\vspace{3mm}
{\small 主仓库冻结版本：7317601b43b914525c21f09e5a60020552ea677b}\par
{\small A17 冻结版本：2ac81b989a2f13d62df9d5a7abd690a33c2189c7}\par
\vspace{7mm}
{\small 交付范围：中文理论报告、独立实验交接文件、推导与微型验证源文件。\par
本轮未启动新的大规模成像或参数扫掠；未修改、提交或推送研究仓库。}
\end{titlepage}
\tableofcontents
\clearpage
$body$
\end{document}
